\section{Complete E4 Operating-Point Cost Accounting}\label{sec:r46-e4-costs}
Tables~\ref{tab:r46-cost-primary-a}--\ref{tab:r46-cost-secondary-b} print all 162 rows of the saved operating-point projection: 81 for the topic-held-out binary study and 81 for the known-topic, template-held-out study. Each study has six raw-reference rows and 15 rows for each of five policies (three local points per target and two imported defaults per target). The two splits can fit different calibrations, so their cost grids are kept separate. Each row uses 500 native-output replay units from 100 topics, with five templates per topic; these are not 500 independent topics. These tables reformat saved results without new captures, attacker fits or statistical tests. Full-precision values, per-attacker scores and pairwise contrasts remain in the recovery-display CSV/JSON outputs.

\paragraph{Keys and provenance}
Q8, Q0.6 and L1 denote qwen3:8b, qwen3:0.6b and llama3.2:1b. Each source/target pair names defense calibration and output generator. Origin B uses the original/continuation recovery mapping and its base raw reference; H uses separately collected hard-cap captures and that study's independent raw reference. A star marks the Q8 fixed-cell point-2 row routed to the separate post hoc 500-session arm. The original run remains incomplete. Row IDs P01--P81 and S01--S81 join the two cost panels within each study.

The zero-based point IDs preserve the machine output: BU uses bucket multipliers 0.5, 1 and 2 of the source-training frame-size $p95$, rounded upward to a 16-byte quantum; BA groups 2, 4 or 8 native chunks; FR emits cells on 20, 50 or 100~ms ticks with source-training $p95$ cell size rounded upward to 16 bytes. FR retains a variable stopping time. MP is minimum event-count padding at source-training quantiles 0.90, 0.95 and 0.99, retaining overflow; HC is the separate hard application-event cap, which is not a cap on TLS record fragments. Imports use point 1. Raw is the corresponding unshaped reference. Identical default BU/BA/FR source parameters yield no calibration-transfer test.

\paragraph{Measured quantities and units}
Panel A gives summed response application TLS bytes $B$, the paired raw sum $B_0$, their ratio, receiver completion percentiles $C_{50},C_{95}$, and percentiles $\Delta C_{50},\Delta C_{95}$ of per-session completion differences against paired raw. Byte counts include TLS record headers and ciphertext overhead; request and handshake traffic is excluded. Completion is elapsed from the replay epoch immediately after the request carrier is received, with the saved schedule aligned to the first native chunk. It is not online model time-to-first-token. A difference of separately computed completion percentiles is not the paired-difference statistic. Negative differences can reflect timing noise across separate captures.

Panel B gives the receiver first-to-last-record span $S_{95}$, original content-chunk endpoint count $N_e$, endpoint-delay percentiles $\Delta R_{50},\Delta R_{95}$ against the same byte endpoints in paired raw captures, planned queue-tail percentiles $T_{50},T_{95}$, and summed sender blocking $W$. All times are milliseconds except $W$, which is seconds. Endpoint percentiles pool chunk endpoints, whereas completion percentiles use sessions. Planned tails are schedule quantities, not measured latency. Summed blocking is not end-to-end runtime. Decimal values are rounded for display; the machine CSV retains precision. Local replay timing does not establish online defended-generation, cloud-service or WAN cost, and unchanged decoded content does not establish online task utility.

\begingroup\fontsize{7.5}{9}\selectfont\setlength{\tabcolsep}{3pt}
\begin{longtable}{llllrrrrrrr}
\caption{Complete primary E4 cost grid, panel A: response byte totals and receiver completion (ms). Keys, units and recovery limits are defined above; all rows retain their own raw reference.}\label{tab:r46-cost-primary-a}\\
\toprule
ID & Origin & Source/target & Policy/point & $B$ & $B_0$ & $B/B_0$ & $C_{50}$ & $C_{95}$ & $\Delta C_{50}$ & $\Delta C_{95}$\\\midrule
\endfirsthead
\multicolumn{11}{l}{\tablename\ \thetable\ continued}\\\toprule
ID & Origin & Source/target & Policy/point & $B$ & $B_0$ & $B/B_0$ & $C_{50}$ & $C_{95}$ & $\Delta C_{50}$ & $\Delta C_{95}$\\\midrule
\endhead
\midrule
\multicolumn{11}{r}{Continued on next page}\\
\endfoot
\bottomrule\endlastfoot
P01 & B & Q8/Q8 & Raw/0 & 1234823 & 1234823 & 1.0000 & 1295.423 & 1829.320 & 0.000 & 0.000\\
P02 & B & Q0.6/Q0.6 & Raw/0 & 1137455 & 1137455 & 1.0000 & 320.535 & 615.328 & 0.000 & 0.000\\
P03 & B & L1/L1 & Raw/0 & 1817075 & 1817075 & 1.0000 & 692.232 & 908.777 & 0.000 & 0.000\\
P04 & B & Q8/Q8 & BU/0 & 1450536 & 1234823 & 1.1747 & 1293.777 & 1830.308 & 0.008 & 1.893\\
P05 & B & Q8/Q8 & BU/1 & 1804952 & 1234823 & 1.4617 & 1295.587 & 1829.407 & 0.023 & 2.269\\
P06 & B & Q8/Q8 & BU/2 & 2335160 & 1234823 & 1.8911 & 1295.075 & 1829.235 & 0.014 & 2.327\\
P07 & B & Q8/Q8 & BA/0 & 725059 & 1234823 & 0.5872 & 1294.610 & 1829.882 & 0.148 & 2.194\\
P08 & B & Q8/Q8 & BA/1 & 470239 & 1234823 & 0.3808 & 1294.175 & 1830.504 & 0.152 & 2.381\\
P09 & B & Q8/Q8 & BA/2 & 342922 & 1234823 & 0.2777 & 1295.306 & 1829.428 & 0.256 & 2.397\\
P10 & B & Q8/Q8 & FR/0 & 1942056 & 1234823 & 1.5727 & 1363.322 & 1903.405 & 87.626 & 211.617\\
P11 & B & Q8/Q8 & FR/1 & 1823040 & 1234823 & 1.4764 & 3250.907 & 4351.173 & 1957.286 & 2594.619\\
P12\textsuperscript{*} & B & Q8/Q8 & FR/2 & 1817694 & 1234823 & 1.4720 & 6500.899 & 8700.875 & 5181.888 & 6896.897\\
P13 & B & Q8/Q8 & MP/0 & 1476840 & 1234823 & 1.1960 & 1580.039 & 1926.689 & 337.938 & 578.010\\
P14 & B & Q8/Q8 & MP/1 & 1560571 & 1234823 & 1.2638 & 1693.060 & 2037.512 & 443.552 & 683.964\\
P15 & B & Q8/Q8 & MP/2 & 1693282 & 1234823 & 1.3713 & 1860.572 & 2199.910 & 620.893 & 860.122\\
P16 & B & Q0.6/Q8 & BU/1 & 1804952 & 1234823 & 1.4617 & 1295.587 & 1829.407 & 0.023 & 2.269\\
P17 & B & Q0.6/Q8 & BA/1 & 470239 & 1234823 & 0.3808 & 1294.175 & 1830.504 & 0.152 & 2.381\\
P18 & B & Q0.6/Q8 & FR/1 & 1823040 & 1234823 & 1.4764 & 3250.907 & 4351.173 & 1957.286 & 2594.619\\
P19 & B & Q0.6/Q8 & MP/1 & 1572072 & 1234823 & 1.2731 & 1722.010 & 2061.848 & 460.678 & 701.282\\
P20 & B & L1/Q8 & BU/1 & 1804952 & 1234823 & 1.4617 & 1295.587 & 1829.407 & 0.023 & 2.269\\
P21 & B & L1/Q8 & BA/1 & 470239 & 1234823 & 0.3808 & 1294.175 & 1830.504 & 0.152 & 2.381\\
P22 & B & L1/Q8 & FR/1 & 1823040 & 1234823 & 1.4764 & 3250.907 & 4351.173 & 1957.286 & 2594.619\\
P23 & B & L1/Q8 & MP/1 & 2191545 & 1234823 & 1.7748 & 2501.756 & 2839.471 & 1261.563 & 1501.588\\
P24 & B & Q8/Q0.6 & BU/1 & 1667306 & 1137455 & 1.4658 & 320.270 & 618.676 & 0.027 & 2.147\\
P25 & B & Q8/Q0.6 & BA/1 & 432236 & 1137455 & 0.3800 & 321.274 & 615.270 & 0.243 & 2.217\\
P26 & B & Q8/Q0.6 & FR/1 & 1679400 & 1137455 & 1.4765 & 2950.815 & 4351.674 & 2598.551 & 3916.426\\
P27 & B & Q8/Q0.6 & MP/1 & 1546562 & 1137455 & 1.3597 & 892.849 & 1218.958 & 578.112 & 943.546\\
P28 & B & Q0.6/Q0.6 & BU/0 & 1338538 & 1137455 & 1.1768 & 321.288 & 615.274 & 0.016 & 1.997\\
P29 & B & Q0.6/Q0.6 & BU/1 & 1667306 & 1137455 & 1.4658 & 320.270 & 618.676 & 0.027 & 2.147\\
P30 & B & Q0.6/Q0.6 & BU/2 & 2156730 & 1137455 & 1.8961 & 321.299 & 615.208 & -0.013 & 2.206\\
P31 & B & Q0.6/Q0.6 & BA/0 & 667309 & 1137455 & 0.5867 & 320.641 & 615.441 & 0.086 & 2.179\\
P32 & B & Q0.6/Q0.6 & BA/1 & 432236 & 1137455 & 0.3800 & 321.274 & 615.270 & 0.243 & 2.217\\
P33 & B & Q0.6/Q0.6 & BA/2 & 314560 & 1137455 & 0.2765 & 321.129 & 616.094 & 0.357 & 2.991\\
P34 & B & Q0.6/Q0.6 & FR/0 & 1687770 & 1137455 & 1.4838 & 1180.883 & 1742.329 & 855.608 & 1321.031\\
P35 & B & Q0.6/Q0.6 & FR/1 & 1679400 & 1137455 & 1.4765 & 2950.815 & 4351.674 & 2598.551 & 3916.426\\
P36 & B & Q0.6/Q0.6 & FR/2 & 1678860 & 1137455 & 1.4760 & 5900.661 & 8702.107 & 5493.071 & 8328.051\\
P37 & B & Q0.6/Q0.6 & MP/0 & 1483818 & 1137455 & 1.3045 & 806.204 & 1137.179 & 480.457 & 859.217\\
P38 & B & Q0.6/Q0.6 & MP/1 & 1558838 & 1137455 & 1.3705 & 908.853 & 1247.212 & 585.866 & 960.802\\
P39 & B & Q0.6/Q0.6 & MP/2 & 1803273 & 1137455 & 1.5854 & 1225.973 & 1565.247 & 919.783 & 1280.435\\
P40 & B & L1/Q0.6 & BU/1 & 1667306 & 1137455 & 1.4658 & 320.270 & 618.676 & 0.027 & 2.147\\
P41 & B & L1/Q0.6 & BA/1 & 432236 & 1137455 & 0.3800 & 321.274 & 615.270 & 0.243 & 2.217\\
P42 & B & L1/Q0.6 & FR/1 & 1679400 & 1137455 & 1.4765 & 2950.815 & 4351.674 & 2598.551 & 3916.426\\
P43 & B & L1/Q0.6 & MP/1 & 2173196 & 1137455 & 1.9106 & 1712.972 & 2043.229 & 1382.089 & 1760.152\\
P44 & B & Q8/L1 & BU/1 & 2675540 & 1817075 & 1.4724 & 692.325 & 908.943 & -0.036 & 1.954\\
P45 & B & Q8/L1 & BA/1 & 677019 & 1817075 & 0.3726 & 693.477 & 909.000 & 0.269 & 2.474\\
P46 & B & Q8/L1 & FR/1 & 2687094 & 1817075 & 1.4788 & 4803.182 & 6401.016 & 4133.387 & 5495.639\\
P47 & B & Q8/L1 & MP/1 & 1854244 & 1817075 & 1.0205 & 742.145 & 911.919 & 0.327 & 280.523\\
P48 & B & Q0.6/L1 & BU/1 & 2675540 & 1817075 & 1.4724 & 692.325 & 908.943 & -0.036 & 1.954\\
P49 & B & Q0.6/L1 & BA/1 & 677019 & 1817075 & 0.3726 & 693.477 & 909.000 & 0.269 & 2.474\\
P50 & B & Q0.6/L1 & FR/1 & 2687094 & 1817075 & 1.4788 & 4803.182 & 6401.016 & 4133.387 & 5495.639\\
P51 & B & Q0.6/L1 & MP/1 & 1857871 & 1817075 & 1.0225 & 748.050 & 924.939 & 0.519 & 300.033\\
P52 & B & L1/L1 & BU/0 & 2132516 & 1817075 & 1.1736 & 693.267 & 908.796 & 0.048 & 2.280\\
P53 & B & L1/L1 & BU/1 & 2675540 & 1817075 & 1.4724 & 692.325 & 908.943 & -0.036 & 1.954\\
P54 & B & L1/L1 & BU/2 & 3463700 & 1817075 & 1.9062 & 693.183 & 908.590 & -0.032 & 2.536\\
P55 & B & L1/L1 & BA/0 & 1056986 & 1817075 & 0.5817 & 693.753 & 908.383 & 0.135 & 2.534\\
P56 & B & L1/L1 & BA/1 & 677019 & 1817075 & 0.3726 & 693.477 & 909.000 & 0.269 & 2.474\\
P57 & B & L1/L1 & BA/2 & 487299 & 1817075 & 0.2682 & 693.224 & 909.089 & 0.332 & 2.642\\
P58 & B & L1/L1 & FR/0 & 2688768 & 1817075 & 1.4797 & 1922.339 & 2560.925 & 1244.967 & 1655.937\\
P59 & B & L1/L1 & FR/1 & 2687094 & 1817075 & 1.4788 & 4803.182 & 6401.016 & 4133.387 & 5495.639\\
P60 & B & L1/L1 & FR/2 & 2687040 & 1817075 & 1.4788 & 9603.091 & 12801.128 & 8982.591 & 11895.140\\
P61 & B & L1/L1 & MP/0 & 2248750 & 1817075 & 1.2376 & 1284.654 & 1601.880 & 599.147 & 1059.158\\
P62 & B & L1/L1 & MP/1 & 2274015 & 1817075 & 1.2515 & 1318.023 & 1637.724 & 639.696 & 1096.149\\
P63 & B & L1/L1 & MP/2 & 2274015 & 1817075 & 1.2515 & 1318.023 & 1637.724 & 639.696 & 1096.149\\
P64 & H & Q8/Q8 & Raw/0 & 1234823 & 1234823 & 1.0000 & 1294.373 & 1831.319 & 0.000 & 0.000\\
P65 & H & Q0.6/Q0.6 & Raw/0 & 1137455 & 1137455 & 1.0000 & 320.551 & 615.135 & 0.000 & 0.000\\
P66 & H & L1/L1 & Raw/0 & 1817075 & 1817075 & 1.0000 & 693.457 & 908.990 & 0.000 & 0.000\\
P67 & H & Q8/Q8 & HC/0 & 1466145 & 1234823 & 1.1873 & 1580.311 & 1928.678 & 336.041 & 579.044\\
P68 & H & Q8/Q8 & HC/1 & 1556045 & 1234823 & 1.2601 & 1692.538 & 2037.463 & 442.281 & 683.558\\
P69 & H & Q8/Q8 & HC/2 & 1692445 & 1234823 & 1.3706 & 1860.567 & 2201.509 & 620.735 & 860.399\\
P70 & H & Q0.6/Q8 & HC/1 & 1568445 & 1234823 & 1.2702 & 1721.794 & 2066.016 & 460.672 & 701.724\\
P71 & H & L1/Q8 & HC/1 & 2191545 & 1234823 & 1.7748 & 2502.074 & 2838.978 & 1261.416 & 1503.480\\
P72 & H & Q8/Q0.6 & HC/1 & 1537696 & 1137455 & 1.3519 & 894.208 & 1218.811 & 577.889 & 943.325\\
P73 & H & Q0.6/Q0.6 & HC/0 & 1469496 & 1137455 & 1.2919 & 806.650 & 1136.270 & 480.121 & 859.207\\
P74 & H & Q0.6/Q0.6 & HC/1 & 1550096 & 1137455 & 1.3628 & 906.981 & 1246.250 & 581.941 & 960.336\\
P75 & H & Q0.6/Q0.6 & HC/2 & 1801196 & 1137455 & 1.5835 & 1225.473 & 1567.093 & 919.592 & 1281.846\\
P76 & H & L1/Q0.6 & HC/1 & 2173196 & 1137455 & 1.9106 & 1712.330 & 2044.914 & 1383.144 & 1759.913\\
P77 & H & Q8/L1 & HC/1 & 1638515 & 1817075 & 0.9017 & 740.999 & 911.919 & 0.937 & 280.521\\
P78 & H & Q0.6/L1 & HC/1 & 1650915 & 1817075 & 0.9086 & 748.201 & 925.883 & 1.068 & 295.629\\
P79 & H & L1/L1 & HC/0 & 2246115 & 1817075 & 1.2361 & 1284.388 & 1601.052 & 598.696 & 1060.263\\
P80 & H & L1/L1 & HC/1 & 2274015 & 1817075 & 1.2515 & 1318.254 & 1638.009 & 638.519 & 1099.758\\
P81 & H & L1/L1 & HC/2 & 2274015 & 1817075 & 1.2515 & 1318.254 & 1638.009 & 638.519 & 1099.758\\
\end{longtable}\endgroup

\begingroup\fontsize{7.5}{9}\selectfont\setlength{\tabcolsep}{3pt}
\begin{longtable}{lrrrrrrr}
\caption{Complete primary E4 cost grid, panel B: receiver span, chunk delivery, planned tail (ms), and sender blocking (s). Keys, units and recovery limits are defined above; all rows retain their own raw reference.}\label{tab:r46-cost-primary-b}\\
\toprule
ID & $S_{95}$ & $N_e$ & $\Delta R_{50}$ & $\Delta R_{95}$ & $T_{50}$ & $T_{95}$ & $W$ (s)\\\midrule
\endfirsthead
\multicolumn{8}{l}{\tablename\ \thetable\ continued}\\\toprule
ID & $S_{95}$ & $N_e$ & $\Delta R_{50}$ & $\Delta R_{95}$ & $T_{50}$ & $T_{95}$ & $W$ (s)\\\midrule
\endhead
\midrule
\multicolumn{8}{r}{Continued on next page}\\
\endfoot
\bottomrule\endlastfoot
P01 & 1829.117 & 33138 & 0.000 & 0.000 & 0.000 & 0.000 & 9.115\\
P02 & 615.067 & 30589 & 0.000 & 0.000 & 0.000 & 0.000 & 7.739\\
P03 & 908.490 & 49260 & 0.000 & 0.000 & 0.000 & 0.000 & 12.500\\
P04 & 1829.981 & 33138 & -0.011 & 2.118 & 0.000 & 0.000 & 8.879\\
P05 & 1828.941 & 33138 & -0.008 & 2.166 & 0.000 & 0.000 & 8.888\\
P06 & 1829.118 & 33138 & -0.018 & 2.039 & 0.000 & 0.000 & 8.983\\
P07 & 1786.747 & 33138 & 3.354 & 23.748 & 0.000 & 0.000 & 4.494\\
P08 & 1733.796 & 33138 & 27.591 & 65.730 & 0.000 & 0.000 & 2.453\\
P09 & 1651.834 & 33138 & 66.356 & 141.359 & 0.000 & 0.000 & 1.273\\
P10 & 1902.783 & 33138 & 38.348 & 126.055 & 87.589 & 211.466 & 9.771\\
P11 & 4350.777 & 33138 & 971.764 & 2052.567 & 1957.568 & 2594.453 & 9.122\\
P12\textsuperscript{*} & 8700.601 & 33138 & 2599.329 & 5430.918 & 5183.798 & 6898.355 & 10.646\\
P13 & 1926.345 & 33138 & -0.008 & 2.121 & 340.000 & 580.000 & 11.093\\
P14 & 2037.094 & 33138 & -0.017 & 2.209 & 440.000 & 681.000 & 11.788\\
P15 & 2199.573 & 33138 & -0.004 & 2.152 & 620.000 & 860.000 & 12.918\\
P16 & 1828.941 & 33138 & -0.008 & 2.166 & 0.000 & 0.000 & 8.888\\
P17 & 1733.796 & 33138 & 27.591 & 65.730 & 0.000 & 0.000 & 2.453\\
P18 & 4350.777 & 33138 & 971.764 & 2052.567 & 1957.568 & 2594.453 & 9.122\\
P19 & 2061.589 & 33138 & -0.014 & 2.282 & 460.000 & 700.000 & 12.083\\
P20 & 1828.941 & 33138 & -0.008 & 2.166 & 0.000 & 0.000 & 8.888\\
P21 & 1733.796 & 33138 & 27.591 & 65.730 & 0.000 & 0.000 & 2.453\\
P22 & 4350.777 & 33138 & 971.764 & 2052.567 & 1957.568 & 2594.453 & 9.122\\
P23 & 2839.104 & 33138 & -0.044 & 2.026 & 1260.000 & 1500.000 & 17.591\\
P24 & 618.232 & 30589 & 0.014 & 1.964 & 0.000 & 0.000 & 7.235\\
P25 & 576.648 & 30589 & 7.508 & 22.289 & 0.000 & 0.000 & 2.276\\
P26 & 4351.429 & 30589 & 1328.791 & 3097.641 & 2597.828 & 3915.933 & 8.327\\
P27 & 1218.567 & 30589 & -0.002 & 1.963 & 580.000 & 941.000 & 10.710\\
P28 & 615.023 & 30589 & 0.014 & 1.956 & 0.000 & 0.000 & 7.631\\
P29 & 618.232 & 30589 & 0.014 & 1.964 & 0.000 & 0.000 & 7.235\\
P30 & 615.011 & 30589 & -0.003 & 1.954 & 0.000 & 0.000 & 7.397\\
P31 & 589.996 & 30589 & 2.256 & 8.854 & 0.000 & 0.000 & 4.139\\
P32 & 576.648 & 30589 & 7.508 & 22.289 & 0.000 & 0.000 & 2.276\\
P33 & 553.723 & 30589 & 17.497 & 47.719 & 0.000 & 0.000 & 1.155\\
P34 & 1742.112 & 30589 & 425.312 & 1017.965 & 855.923 & 1321.829 & 8.411\\
P35 & 4351.429 & 30589 & 1328.791 & 3097.641 & 2597.828 & 3915.933 & 8.327\\
P36 & 8701.612 & 30589 & 2832.336 & 6569.634 & 5493.302 & 8327.522 & 8.415\\
P37 & 1136.720 & 30589 & 0.016 & 2.181 & 480.000 & 860.000 & 10.988\\
P38 & 1246.932 & 30589 & 0.020 & 2.047 & 580.000 & 960.000 & 10.924\\
P39 & 1564.854 & 30589 & 0.004 & 1.974 & 920.000 & 1280.000 & 13.085\\
P40 & 618.232 & 30589 & 0.014 & 1.964 & 0.000 & 0.000 & 7.235\\
P41 & 576.648 & 30589 & 7.508 & 22.289 & 0.000 & 0.000 & 2.276\\
P42 & 4351.429 & 30589 & 1328.791 & 3097.641 & 2597.828 & 3915.933 & 8.327\\
P43 & 2043.095 & 30589 & 0.006 & 1.925 & 1380.000 & 1760.000 & 16.387\\
P44 & 908.448 & 49260 & 0.006 & 2.092 & 0.000 & 0.000 & 12.582\\
P45 & 881.325 & 49260 & 10.607 & 23.258 & 0.000 & 0.000 & 3.409\\
P46 & 6400.732 & 49260 & 2097.878 & 4430.952 & 4130.402 & 5495.072 & 13.908\\
P47 & 911.535 & 49260 & -0.015 & 2.012 & 0.000 & 280.000 & 13.905\\
P48 & 908.448 & 49260 & 0.006 & 2.092 & 0.000 & 0.000 & 12.582\\
P49 & 881.325 & 49260 & 10.607 & 23.258 & 0.000 & 0.000 & 3.409\\
P50 & 6400.732 & 49260 & 2097.878 & 4430.952 & 4130.402 & 5495.072 & 13.908\\
P51 & 924.638 & 49260 & 0.016 & 2.090 & 0.000 & 300.000 & 13.634\\
P52 & 908.417 & 49260 & 0.001 & 2.219 & 0.000 & 0.000 & 12.892\\
P53 & 908.448 & 49260 & 0.006 & 2.092 & 0.000 & 0.000 & 12.582\\
P54 & 908.183 & 49260 & 0.013 & 2.064 & 0.000 & 0.000 & 12.502\\
P55 & 896.956 & 49260 & 2.989 & 9.118 & 0.000 & 0.000 & 6.977\\
P56 & 881.325 & 49260 & 10.607 & 23.258 & 0.000 & 0.000 & 3.409\\
P57 & 854.283 & 49260 & 24.017 & 50.806 & 0.000 & 0.000 & 1.775\\
P58 & 2560.610 & 49260 & 626.775 & 1336.190 & 1244.206 & 1657.096 & 13.470\\
P59 & 6400.732 & 49260 & 2097.878 & 4430.952 & 4130.402 & 5495.072 & 13.908\\
P60 & 12800.843 & 49260 & 4550.378 & 9601.888 & 8982.094 & 11894.334 & 13.146\\
P61 & 1601.507 & 49260 & 0.013 & 2.009 & 600.000 & 1060.000 & 16.316\\
P62 & 1637.454 & 49260 & -0.019 & 1.842 & 640.000 & 1100.000 & 16.607\\
P63 & 1637.454 & 49260 & -0.019 & 1.842 & 640.000 & 1100.000 & 16.607\\
P64 & 1831.069 & 33138 & 0.000 & 0.000 & 0.000 & 0.000 & 8.452\\
P65 & 615.033 & 30589 & 0.000 & 0.000 & 0.000 & 0.000 & 6.988\\
P66 & 908.584 & 49260 & 0.000 & 0.000 & 0.000 & 0.000 & 11.744\\
P67 & 1928.371 & 33138 & 0.039 & 2.560 & 340.000 & 580.000 & 10.273\\
P68 & 2037.127 & 33138 & -0.002 & 2.549 & 440.000 & 681.000 & 10.667\\
P69 & 2201.181 & 33138 & -0.017 & 2.381 & 620.000 & 860.000 & 11.695\\
P70 & 2065.790 & 33138 & 0.015 & 2.543 & 460.000 & 700.000 & 10.927\\
P71 & 2838.726 & 33138 & 0.043 & 2.489 & 1260.000 & 1500.000 & 16.327\\
P72 & 1218.624 & 30589 & -0.014 & 2.173 & 580.000 & 941.000 & 9.901\\
P73 & 1136.157 & 30589 & 0.004 & 2.552 & 480.000 & 860.000 & 9.423\\
P74 & 1246.067 & 30589 & -0.004 & 2.252 & 580.000 & 960.000 & 10.028\\
P75 & 1566.726 & 30589 & -0.018 & 2.088 & 920.000 & 1280.000 & 12.147\\
P76 & 2044.521 & 30589 & -0.013 & 2.116 & 1380.000 & 1760.000 & 14.882\\
P77 & 911.612 & 49260 & 0.115 & 131.225 & 0.000 & 280.000 & 10.036\\
P78 & 925.665 & 49260 & 0.114 & 126.153 & 0.000 & 300.000 & 10.262\\
P79 & 1600.957 & 49260 & -0.016 & 2.340 & 600.000 & 1060.000 & 14.890\\
P80 & 1637.710 & 49260 & -0.016 & 2.267 & 640.000 & 1100.000 & 15.273\\
P81 & 1637.710 & 49260 & -0.016 & 2.267 & 640.000 & 1100.000 & 15.273\\
\end{longtable}\endgroup

\begingroup\fontsize{7.5}{9}\selectfont\setlength{\tabcolsep}{3pt}
\begin{longtable}{llllrrrrrrr}
\caption{Complete secondary E4 cost grid, panel A: response byte totals and receiver completion (ms). Keys, units and recovery limits are defined above; all rows retain their own raw reference.}\label{tab:r46-cost-secondary-a}\\
\toprule
ID & Origin & Source/target & Policy/point & $B$ & $B_0$ & $B/B_0$ & $C_{50}$ & $C_{95}$ & $\Delta C_{50}$ & $\Delta C_{95}$\\\midrule
\endfirsthead
\multicolumn{11}{l}{\tablename\ \thetable\ continued}\\\toprule
ID & Origin & Source/target & Policy/point & $B$ & $B_0$ & $B/B_0$ & $C_{50}$ & $C_{95}$ & $\Delta C_{50}$ & $\Delta C_{95}$\\\midrule
\endhead
\midrule
\multicolumn{11}{r}{Continued on next page}\\
\endfoot
\bottomrule\endlastfoot
S01 & B & Q8/Q8 & Raw/0 & 1234823 & 1234823 & 1.0000 & 1295.423 & 1829.320 & 0.000 & 0.000\\
S02 & B & Q0.6/Q0.6 & Raw/0 & 1137455 & 1137455 & 1.0000 & 320.535 & 615.328 & 0.000 & 0.000\\
S03 & B & L1/L1 & Raw/0 & 1817075 & 1817075 & 1.0000 & 692.232 & 908.777 & 0.000 & 0.000\\
S04 & B & Q8/Q8 & BU/0 & 1450536 & 1234823 & 1.1747 & 1293.777 & 1830.308 & 0.008 & 1.893\\
S05 & B & Q8/Q8 & BU/1 & 1804952 & 1234823 & 1.4617 & 1295.587 & 1829.407 & 0.023 & 2.269\\
S06 & B & Q8/Q8 & BU/2 & 2335160 & 1234823 & 1.8911 & 1295.075 & 1829.235 & 0.014 & 2.327\\
S07 & B & Q8/Q8 & BA/0 & 725059 & 1234823 & 0.5872 & 1294.610 & 1829.882 & 0.148 & 2.194\\
S08 & B & Q8/Q8 & BA/1 & 470239 & 1234823 & 0.3808 & 1294.175 & 1830.504 & 0.152 & 2.381\\
S09 & B & Q8/Q8 & BA/2 & 342922 & 1234823 & 0.2777 & 1295.306 & 1829.428 & 0.256 & 2.397\\
S10 & B & Q8/Q8 & FR/0 & 1942056 & 1234823 & 1.5727 & 1363.322 & 1903.405 & 87.626 & 211.617\\
S11 & B & Q8/Q8 & FR/1 & 1823040 & 1234823 & 1.4764 & 3250.907 & 4351.173 & 1957.286 & 2594.619\\
S12\textsuperscript{*} & B & Q8/Q8 & FR/2 & 1817694 & 1234823 & 1.4720 & 6500.899 & 8700.875 & 5181.888 & 6896.897\\
S13 & B & Q8/Q8 & MP/0 & 1478359 & 1234823 & 1.1972 & 1599.728 & 1944.713 & 339.744 & 602.360\\
S14 & B & Q8/Q8 & MP/1 & 1555611 & 1234823 & 1.2598 & 1698.198 & 2048.799 & 440.384 & 701.457\\
S15 & B & Q8/Q8 & MP/2 & 1679022 & 1234823 & 1.3597 & 1857.490 & 2193.598 & 601.644 & 862.762\\
S16 & B & Q0.6/Q8 & BU/1 & 1804952 & 1234823 & 1.4617 & 1295.587 & 1829.407 & 0.023 & 2.269\\
S17 & B & Q0.6/Q8 & BA/1 & 470239 & 1234823 & 0.3808 & 1294.175 & 1830.504 & 0.152 & 2.381\\
S18 & B & Q0.6/Q8 & FR/1 & 1823040 & 1234823 & 1.4764 & 3250.907 & 4351.173 & 1957.286 & 2594.619\\
S19 & B & Q0.6/Q8 & MP/1 & 1566709 & 1234823 & 1.2688 & 1739.204 & 2066.977 & 457.443 & 719.974\\
S20 & B & L1/Q8 & BU/1 & 1804952 & 1234823 & 1.4617 & 1295.587 & 1829.407 & 0.023 & 2.269\\
S21 & B & L1/Q8 & BA/1 & 470239 & 1234823 & 0.3808 & 1294.175 & 1830.504 & 0.152 & 2.381\\
S22 & B & L1/Q8 & FR/1 & 1823040 & 1234823 & 1.4764 & 3250.907 & 4351.173 & 1957.286 & 2594.619\\
S23 & B & L1/Q8 & MP/1 & 2191545 & 1234823 & 1.7748 & 2501.756 & 2839.471 & 1261.563 & 1501.588\\
S24 & B & Q8/Q0.6 & BU/1 & 1667306 & 1137455 & 1.4658 & 320.270 & 618.676 & 0.027 & 2.147\\
S25 & B & Q8/Q0.6 & BA/1 & 432236 & 1137455 & 0.3800 & 321.274 & 615.270 & 0.243 & 2.217\\
S26 & B & Q8/Q0.6 & FR/1 & 1679400 & 1137455 & 1.4765 & 2950.815 & 4351.674 & 2598.551 & 3916.426\\
S27 & B & Q8/Q0.6 & MP/1 & 1538719 & 1137455 & 1.3528 & 880.137 & 1234.519 & 559.818 & 960.371\\
S28 & B & Q0.6/Q0.6 & BU/0 & 1338538 & 1137455 & 1.1768 & 321.288 & 615.274 & 0.016 & 1.997\\
S29 & B & Q0.6/Q0.6 & BU/1 & 1667306 & 1137455 & 1.4658 & 320.270 & 618.676 & 0.027 & 2.147\\
S30 & B & Q0.6/Q0.6 & BU/2 & 2156730 & 1137455 & 1.8961 & 321.299 & 615.208 & -0.013 & 2.206\\
S31 & B & Q0.6/Q0.6 & BA/0 & 667309 & 1137455 & 0.5867 & 320.641 & 615.441 & 0.086 & 2.179\\
S32 & B & Q0.6/Q0.6 & BA/1 & 432236 & 1137455 & 0.3800 & 321.274 & 615.270 & 0.243 & 2.217\\
S33 & B & Q0.6/Q0.6 & BA/2 & 314560 & 1137455 & 0.2765 & 321.129 & 616.094 & 0.357 & 2.991\\
S34 & B & Q0.6/Q0.6 & FR/0 & 1687770 & 1137455 & 1.4838 & 1180.883 & 1742.329 & 855.608 & 1321.031\\
S35 & B & Q0.6/Q0.6 & FR/1 & 1679400 & 1137455 & 1.4765 & 2950.815 & 4351.674 & 2598.551 & 3916.426\\
S36 & B & Q0.6/Q0.6 & FR/2 & 1678860 & 1137455 & 1.4760 & 5900.661 & 8702.107 & 5493.071 & 8328.051\\
S37 & B & Q0.6/Q0.6 & MP/0 & 1485151 & 1137455 & 1.3057 & 822.327 & 1174.713 & 500.032 & 900.432\\
S38 & B & Q0.6/Q0.6 & MP/1 & 1556420 & 1137455 & 1.3683 & 909.761 & 1285.201 & 599.782 & 1004.551\\
S39 & B & Q0.6/Q0.6 & MP/2 & 1788579 & 1137455 & 1.5724 & 1226.888 & 1583.471 & 918.192 & 1300.856\\
S40 & B & L1/Q0.6 & BU/1 & 1667306 & 1137455 & 1.4658 & 320.270 & 618.676 & 0.027 & 2.147\\
S41 & B & L1/Q0.6 & BA/1 & 432236 & 1137455 & 0.3800 & 321.274 & 615.270 & 0.243 & 2.217\\
S42 & B & L1/Q0.6 & FR/1 & 1679400 & 1137455 & 1.4765 & 2950.815 & 4351.674 & 2598.551 & 3916.426\\
S43 & B & L1/Q0.6 & MP/1 & 2173196 & 1137455 & 1.9106 & 1712.972 & 2043.229 & 1382.089 & 1760.152\\
S44 & B & Q8/L1 & BU/1 & 2675540 & 1817075 & 1.4724 & 692.325 & 908.943 & -0.036 & 1.954\\
S45 & B & Q8/L1 & BA/1 & 677019 & 1817075 & 0.3726 & 693.477 & 909.000 & 0.269 & 2.474\\
S46 & B & Q8/L1 & FR/1 & 2687094 & 1817075 & 1.4788 & 4803.182 & 6401.016 & 4133.387 & 5495.639\\
S47 & B & Q8/L1 & MP/1 & 1856321 & 1817075 & 1.0216 & 745.348 & 924.261 & 0.396 & 280.610\\
S48 & B & Q0.6/L1 & BU/1 & 2675540 & 1817075 & 1.4724 & 692.325 & 908.943 & -0.036 & 1.954\\
S49 & B & Q0.6/L1 & BA/1 & 677019 & 1817075 & 0.3726 & 693.477 & 909.000 & 0.269 & 2.474\\
S50 & B & Q0.6/L1 & FR/1 & 2687094 & 1817075 & 1.4788 & 4803.182 & 6401.016 & 4133.387 & 5495.639\\
S51 & B & Q0.6/L1 & MP/1 & 1860537 & 1817075 & 1.0239 & 751.845 & 929.916 & 0.432 & 320.119\\
S52 & B & L1/L1 & BU/0 & 2132516 & 1817075 & 1.1736 & 693.267 & 908.796 & 0.048 & 2.280\\
S53 & B & L1/L1 & BU/1 & 2675540 & 1817075 & 1.4724 & 692.325 & 908.943 & -0.036 & 1.954\\
S54 & B & L1/L1 & BU/2 & 3463700 & 1817075 & 1.9062 & 693.183 & 908.590 & -0.032 & 2.536\\
S55 & B & L1/L1 & BA/0 & 1056986 & 1817075 & 0.5817 & 693.753 & 908.383 & 0.135 & 2.534\\
S56 & B & L1/L1 & BA/1 & 677019 & 1817075 & 0.3726 & 693.477 & 909.000 & 0.269 & 2.474\\
S57 & B & L1/L1 & BA/2 & 487299 & 1817075 & 0.2682 & 693.224 & 909.089 & 0.332 & 2.642\\
S58 & B & L1/L1 & FR/0 & 2688768 & 1817075 & 1.4797 & 1922.339 & 2560.925 & 1244.967 & 1655.937\\
S59 & B & L1/L1 & FR/1 & 2687094 & 1817075 & 1.4788 & 4803.182 & 6401.016 & 4133.387 & 5495.639\\
S60 & B & L1/L1 & FR/2 & 2687040 & 1817075 & 1.4788 & 9603.091 & 12801.128 & 8982.591 & 11895.140\\
S61 & B & L1/L1 & MP/0 & 2245681 & 1817075 & 1.2359 & 1276.205 & 1630.154 & 580.461 & 1079.246\\
S62 & B & L1/L1 & MP/1 & 2274015 & 1817075 & 1.2515 & 1318.023 & 1637.724 & 639.696 & 1096.149\\
S63 & B & L1/L1 & MP/2 & 2274015 & 1817075 & 1.2515 & 1318.023 & 1637.724 & 639.696 & 1096.149\\
S64 & H & Q8/Q8 & Raw/0 & 1234823 & 1234823 & 1.0000 & 1294.373 & 1831.319 & 0.000 & 0.000\\
S65 & H & Q0.6/Q0.6 & Raw/0 & 1137455 & 1137455 & 1.0000 & 320.551 & 615.135 & 0.000 & 0.000\\
S66 & H & L1/L1 & Raw/0 & 1817075 & 1817075 & 1.0000 & 693.457 & 908.990 & 0.000 & 0.000\\
S67 & H & Q8/Q8 & HC/0 & 1459945 & 1234823 & 1.1823 & 1599.442 & 1944.733 & 339.127 & 600.627\\
S68 & H & Q8/Q8 & HC/1 & 1546745 & 1234823 & 1.2526 & 1699.962 & 2049.898 & 439.988 & 701.765\\
S69 & H & Q8/Q8 & HC/2 & 1676945 & 1234823 & 1.3580 & 1857.273 & 2193.011 & 602.001 & 861.035\\
S70 & H & Q0.6/Q8 & HC/1 & 1562245 & 1234823 & 1.2652 & 1739.884 & 2066.989 & 458.103 & 719.233\\
S71 & H & L1/Q8 & HC/1 & 2191545 & 1234823 & 1.7748 & 2502.074 & 2838.978 & 1261.416 & 1503.480\\
S72 & H & Q8/Q0.6 & HC/1 & 1528396 & 1137455 & 1.3437 & 879.492 & 1235.151 & 560.224 & 960.887\\
S73 & H & Q0.6/Q0.6 & HC/0 & 1457096 & 1137455 & 1.2810 & 823.027 & 1174.287 & 499.948 & 899.929\\
S74 & H & Q0.6/Q0.6 & HC/1 & 1543896 & 1137455 & 1.3573 & 908.562 & 1287.710 & 599.640 & 1004.156\\
S75 & H & Q0.6/Q0.6 & HC/2 & 1785696 & 1137455 & 1.5699 & 1225.799 & 1583.646 & 918.593 & 1300.705\\
S76 & H & L1/Q0.6 & HC/1 & 2173196 & 1137455 & 1.9106 & 1712.330 & 2044.914 & 1383.144 & 1759.913\\
S77 & H & Q8/L1 & HC/1 & 1629215 & 1817075 & 0.8966 & 746.786 & 926.098 & 0.942 & 282.112\\
S78 & H & Q0.6/L1 & HC/1 & 1644715 & 1817075 & 0.9051 & 751.150 & 929.332 & 1.049 & 320.401\\
S79 & H & L1/L1 & HC/0 & 2239915 & 1817075 & 1.2327 & 1277.044 & 1630.239 & 579.650 & 1078.577\\
S80 & H & L1/L1 & HC/1 & 2274015 & 1817075 & 1.2515 & 1318.254 & 1638.009 & 638.519 & 1099.758\\
S81 & H & L1/L1 & HC/2 & 2274015 & 1817075 & 1.2515 & 1318.254 & 1638.009 & 638.519 & 1099.758\\
\end{longtable}\endgroup

\begingroup\fontsize{7.5}{9}\selectfont\setlength{\tabcolsep}{3pt}
\begin{longtable}{lrrrrrrr}
\caption{Complete secondary E4 cost grid, panel B: receiver span, chunk delivery, planned tail (ms), and sender blocking (s). Keys, units and recovery limits are defined above; all rows retain their own raw reference.}\label{tab:r46-cost-secondary-b}\\
\toprule
ID & $S_{95}$ & $N_e$ & $\Delta R_{50}$ & $\Delta R_{95}$ & $T_{50}$ & $T_{95}$ & $W$ (s)\\\midrule
\endfirsthead
\multicolumn{8}{l}{\tablename\ \thetable\ continued}\\\toprule
ID & $S_{95}$ & $N_e$ & $\Delta R_{50}$ & $\Delta R_{95}$ & $T_{50}$ & $T_{95}$ & $W$ (s)\\\midrule
\endhead
\midrule
\multicolumn{8}{r}{Continued on next page}\\
\endfoot
\bottomrule\endlastfoot
S01 & 1829.117 & 33138 & 0.000 & 0.000 & 0.000 & 0.000 & 9.115\\
S02 & 615.067 & 30589 & 0.000 & 0.000 & 0.000 & 0.000 & 7.739\\
S03 & 908.490 & 49260 & 0.000 & 0.000 & 0.000 & 0.000 & 12.500\\
S04 & 1829.981 & 33138 & -0.011 & 2.118 & 0.000 & 0.000 & 8.879\\
S05 & 1828.941 & 33138 & -0.008 & 2.166 & 0.000 & 0.000 & 8.888\\
S06 & 1829.118 & 33138 & -0.018 & 2.039 & 0.000 & 0.000 & 8.983\\
S07 & 1786.747 & 33138 & 3.354 & 23.748 & 0.000 & 0.000 & 4.494\\
S08 & 1733.796 & 33138 & 27.591 & 65.730 & 0.000 & 0.000 & 2.453\\
S09 & 1651.834 & 33138 & 66.356 & 141.359 & 0.000 & 0.000 & 1.273\\
S10 & 1902.783 & 33138 & 38.348 & 126.055 & 87.589 & 211.466 & 9.771\\
S11 & 4350.777 & 33138 & 971.764 & 2052.567 & 1957.568 & 2594.453 & 9.122\\
S12\textsuperscript{*} & 8700.601 & 33138 & 2599.329 & 5430.918 & 5183.798 & 6898.355 & 10.646\\
S13 & 1944.352 & 33138 & 0.000 & 2.220 & 340.000 & 600.000 & 10.976\\
S14 & 2048.434 & 33138 & -0.036 & 1.966 & 440.000 & 700.000 & 11.494\\
S15 & 2193.155 & 33138 & -0.028 & 2.129 & 600.000 & 860.000 & 12.630\\
S16 & 1828.941 & 33138 & -0.008 & 2.166 & 0.000 & 0.000 & 8.888\\
S17 & 1733.796 & 33138 & 27.591 & 65.730 & 0.000 & 0.000 & 2.453\\
S18 & 4350.777 & 33138 & 971.764 & 2052.567 & 1957.568 & 2594.453 & 9.122\\
S19 & 2066.699 & 33138 & -0.015 & 2.100 & 460.000 & 720.000 & 11.887\\
S20 & 1828.941 & 33138 & -0.008 & 2.166 & 0.000 & 0.000 & 8.888\\
S21 & 1733.796 & 33138 & 27.591 & 65.730 & 0.000 & 0.000 & 2.453\\
S22 & 4350.777 & 33138 & 971.764 & 2052.567 & 1957.568 & 2594.453 & 9.122\\
S23 & 2839.104 & 33138 & -0.044 & 2.026 & 1260.000 & 1500.000 & 17.591\\
S24 & 618.232 & 30589 & 0.014 & 1.964 & 0.000 & 0.000 & 7.235\\
S25 & 576.648 & 30589 & 7.508 & 22.289 & 0.000 & 0.000 & 2.276\\
S26 & 4351.429 & 30589 & 1328.791 & 3097.641 & 2597.828 & 3915.933 & 8.327\\
S27 & 1234.409 & 30589 & -0.001 & 1.844 & 560.000 & 960.000 & 10.756\\
S28 & 615.023 & 30589 & 0.014 & 1.956 & 0.000 & 0.000 & 7.631\\
S29 & 618.232 & 30589 & 0.014 & 1.964 & 0.000 & 0.000 & 7.235\\
S30 & 615.011 & 30589 & -0.003 & 1.954 & 0.000 & 0.000 & 7.397\\
S31 & 589.996 & 30589 & 2.256 & 8.854 & 0.000 & 0.000 & 4.139\\
S32 & 576.648 & 30589 & 7.508 & 22.289 & 0.000 & 0.000 & 2.276\\
S33 & 553.723 & 30589 & 17.497 & 47.719 & 0.000 & 0.000 & 1.155\\
S34 & 1742.112 & 30589 & 425.312 & 1017.965 & 855.923 & 1321.829 & 8.411\\
S35 & 4351.429 & 30589 & 1328.791 & 3097.641 & 2597.828 & 3915.933 & 8.327\\
S36 & 8701.612 & 30589 & 2832.336 & 6569.634 & 5493.302 & 8327.522 & 8.415\\
S37 & 1174.381 & 30589 & 0.002 & 2.056 & 500.000 & 900.000 & 10.286\\
S38 & 1284.728 & 30589 & 0.016 & 2.036 & 600.000 & 1001.000 & 11.197\\
S39 & 1583.193 & 30589 & 0.013 & 2.107 & 920.000 & 1300.000 & 13.116\\
S40 & 618.232 & 30589 & 0.014 & 1.964 & 0.000 & 0.000 & 7.235\\
S41 & 576.648 & 30589 & 7.508 & 22.289 & 0.000 & 0.000 & 2.276\\
S42 & 4351.429 & 30589 & 1328.791 & 3097.641 & 2597.828 & 3915.933 & 8.327\\
S43 & 2043.095 & 30589 & 0.006 & 1.925 & 1380.000 & 1760.000 & 16.387\\
S44 & 908.448 & 49260 & 0.006 & 2.092 & 0.000 & 0.000 & 12.582\\
S45 & 881.325 & 49260 & 10.607 & 23.258 & 0.000 & 0.000 & 3.409\\
S46 & 6400.732 & 49260 & 2097.878 & 4430.952 & 4130.402 & 5495.072 & 13.908\\
S47 & 924.035 & 49260 & 0.012 & 2.162 & 0.000 & 280.000 & 13.273\\
S48 & 908.448 & 49260 & 0.006 & 2.092 & 0.000 & 0.000 & 12.582\\
S49 & 881.325 & 49260 & 10.607 & 23.258 & 0.000 & 0.000 & 3.409\\
S50 & 6400.732 & 49260 & 2097.878 & 4430.952 & 4130.402 & 5495.072 & 13.908\\
S51 & 929.602 & 49260 & 0.004 & 2.016 & 0.000 & 320.000 & 13.486\\
S52 & 908.417 & 49260 & 0.001 & 2.219 & 0.000 & 0.000 & 12.892\\
S53 & 908.448 & 49260 & 0.006 & 2.092 & 0.000 & 0.000 & 12.582\\
S54 & 908.183 & 49260 & 0.013 & 2.064 & 0.000 & 0.000 & 12.502\\
S55 & 896.956 & 49260 & 2.989 & 9.118 & 0.000 & 0.000 & 6.977\\
S56 & 881.325 & 49260 & 10.607 & 23.258 & 0.000 & 0.000 & 3.409\\
S57 & 854.283 & 49260 & 24.017 & 50.806 & 0.000 & 0.000 & 1.775\\
S58 & 2560.610 & 49260 & 626.775 & 1336.190 & 1244.206 & 1657.096 & 13.470\\
S59 & 6400.732 & 49260 & 2097.878 & 4430.952 & 4130.402 & 5495.072 & 13.908\\
S60 & 12800.843 & 49260 & 4550.378 & 9601.888 & 8982.094 & 11894.334 & 13.146\\
S61 & 1629.767 & 49260 & 0.006 & 1.992 & 580.000 & 1080.000 & 16.475\\
S62 & 1637.454 & 49260 & -0.019 & 1.842 & 640.000 & 1100.000 & 16.607\\
S63 & 1637.454 & 49260 & -0.019 & 1.842 & 640.000 & 1100.000 & 16.607\\
S64 & 1831.069 & 33138 & 0.000 & 0.000 & 0.000 & 0.000 & 8.452\\
S65 & 615.033 & 30589 & 0.000 & 0.000 & 0.000 & 0.000 & 6.988\\
S66 & 908.584 & 49260 & 0.000 & 0.000 & 0.000 & 0.000 & 11.744\\
S67 & 1944.339 & 33138 & 0.015 & 2.843 & 340.000 & 600.000 & 9.855\\
S68 & 2049.640 & 33138 & 0.024 & 2.675 & 440.000 & 700.000 & 10.676\\
S69 & 2192.849 & 33138 & -0.030 & 2.458 & 600.000 & 860.000 & 11.574\\
S70 & 2066.633 & 33138 & -0.011 & 2.537 & 460.000 & 720.000 & 10.594\\
S71 & 2838.726 & 33138 & 0.043 & 2.489 & 1260.000 & 1500.000 & 16.327\\
S72 & 1234.771 & 30589 & -0.020 & 2.300 & 560.000 & 960.000 & 9.773\\
S73 & 1174.121 & 30589 & 0.004 & 2.890 & 500.000 & 900.000 & 9.333\\
S74 & 1287.379 & 30589 & -0.016 & 2.273 & 600.000 & 1001.000 & 9.619\\
S75 & 1583.518 & 30589 & 0.004 & 2.187 & 920.000 & 1300.000 & 12.141\\
S76 & 2044.521 & 30589 & -0.013 & 2.116 & 1380.000 & 1760.000 & 14.882\\
S77 & 926.026 & 49260 & 0.119 & 142.312 & 0.000 & 280.000 & 10.007\\
S78 & 928.968 & 49260 & 0.126 & 135.138 & 0.000 & 320.000 & 10.382\\
S79 & 1629.658 & 49260 & -0.003 & 2.340 & 580.000 & 1080.000 & 15.273\\
S80 & 1637.710 & 49260 & -0.016 & 2.267 & 640.000 & 1100.000 & 15.273\\
S81 & 1637.710 & 49260 & -0.016 & 2.267 & 640.000 & 1100.000 & 15.273\\
\end{longtable}\endgroup

